
\subsection{The martingale control}
\label{sec:martingale_interpretation_f_alpha_budget}

The previous subsection shows that, on a focus region \(F_L\), log-concavity of the core \(p\)
gives a deterministic bound on \(r(x,z)\), and hence on
\[
    \bar{s}_x(z)=\frac{r(x,z)-1}{\lambda}.
\]
This boundedness is useful for standard concentration arguments. There is also a more intrinsic
reason why the same \(f_\alpha\)-quantity appears in the analysis: it is the natural compensator for
the random multiplicative change of the exponential-weights density.

Consider the actual algorithmic trajectory. Let \(q_u\) be the \cewaalg density before round \(u\),
let \(p_u\) be its core, and let \(r_u\) be the corresponding relative density. Also write
\[
    \E_{u-1}[\cdot]
    \triangleq
    \E[\cdot\mid \cG_{u-1}],
\]
where \(\cG_{u-1}\) is the history before round \(u\). Conditionally on \(\cG_{u-1}\), the objects
\(q_u,p_u,r_u\), and the loss \(\ell_u\), are fixed, while \(A_u\sim p_u\) is random.

For a fixed point \(z\in\K\), the estimated loss used by \scewaalg is
\begin{align}
    \hf_u(z)
    =
    \ell_u(A_u)\bar{s}_{u,A_u}(z)
    =
    \frac{\ell_u(A_u)}{\lambda}
    \bigl(r_u(A_u,z)-1\bigr).
    \label{eq:martingale_hf_recentered}
\end{align}
Define its predictable part by
\begin{align}
    \bar{\ell}_u(z)
    \triangleq
    \E_{u-1}[\hf_u(z)].
    \label{eq:bar_ell_predictable_estimator}
\end{align}
This is the recentered version of the optimistic loss \(\tilde{\ell}_u\). Indeed, since
\[
    \tilde{\ell}_u(z)
    =
    \E_{A_u\sim p_u}[\ell_u(A_u)s_{u,A_u}(z)]
\]
and
\[
    \hf_u(z)
    =
    \ell_u(A_u)(s_{u,A_u}(z)-1),
\]
we have
\begin{align}
    \bar{\ell}_u(z)
    =
    \tilde{\ell}_u(z)-\langle \ell_u,p_u\rangle .
    \label{eq:bar_ell_tilde_ell_relation}
\end{align}
The difference between \(\bar{\ell}_u\) and \(\tilde{\ell}_u\) is a scalar independent of \(z\), so it
does not affect the normalized exponential-weights update.

Now decompose the exponential weight at \(z\) as
\begin{align}
    \exp\left(
    -\eta\sum_{u=1}^s \hf_u(z)
    \right)
    =
    \underbrace{
        \exp\left(
        -\eta\sum_{u=1}^s \bar{\ell}_u(z)
        \right)
    }_{\text{predictable part}}
    \underbrace{
        \exp\left(
        -\eta\sum_{u=1}^s
        \bigl(\hf_u(z)-\bar{\ell}_u(z)\bigr)
        \right)
    }_{\text{martingale multiplier}}.
    \label{eq:martingale_weight_decomposition_bar_ell}
\end{align}
The first factor is what the update would look like if we replaced the random one-point estimator by
its conditional mean. The second factor is the random error. The dangerous case is when
\[
    \sum_{u=1}^s
    \bigl(\hf_u(z)-\bar{\ell}_u(z)\bigr)
\]
is negative, because then the exponential weight at \(z\) is multiplied upward.

Let
\begin{align}
    \alpha=\frac{\eta}{\lambda},
    \qquad
    f_\alpha(r)=e^{-\alpha(r-1)}+\alpha(r-1)-1.
\end{align}
The next proposition says that the upward fluctuation in
\eqref{eq:martingale_weight_decomposition_bar_ell} is paid for by the same \(f_\alpha\)-budget that
appears in the variance bound.

\begin{proposition}[Pointwise martingale compensator]
    \label{prop:pointwise_martingale_compensator_bar_ell}
    Fix \(z\in\K\), and define
    \begin{align}
        \Gamma_s(z)
        \triangleq
        \sum_{u=1}^s
        \E_{u-1}\left[
            f_\alpha(r_u(A_u,z))
            \right].
        \label{eq:Gamma_pointwise_f_alpha_budget}
    \end{align}
    Then
    \begin{align}
        M_s(z)
        \triangleq
        \exp\left(
        -\eta\sum_{u=1}^s
        \bigl(\hf_u(z)-\bar{\ell}_u(z)\bigr)
        -
        \Gamma_s(z)
        \right)
        \label{eq:pointwise_f_alpha_supermartingale}
    \end{align}
    is a non-negative supermartingale. Consequently, for every \(\delta\in(0,1)\), with probability
    at least \(1-\delta\),
    \begin{align}
        -\eta\sum_{u=1}^s
        \bigl(\hf_u(z)-\bar{\ell}_u(z)\bigr)
        \le
        \Gamma_s(z)+\log\frac1\delta
        \label{eq:pointwise_exponent_bound_no_alpha_loss}
    \end{align}
    simultaneously for all \(s\ge1\).
\end{proposition}

\begin{proof}
    Fix \(u\) and condition on \(\cG_{u-1}\). Write
    \[
        \rho=r_u(A_u,z),
        \qquad
        W=\ell_u(A_u)(\rho-1).
    \]
    Then
    \[
        \hf_u(z)=\frac{1}{\lambda}W,
        \qquad
        \bar{\ell}_u(z)=\frac{1}{\lambda}\E_{u-1}[W].
    \]
    Hence
    \[
        \eta\bigl(\hf_u(z)-\bar{\ell}_u(z)\bigr)
        =
        \alpha\bigl(W-\E_{u-1}[W]\bigr).
    \]
    Let
    \[
        \phi_\alpha(v)=e^{-\alpha v}+\alpha v-1.
    \]
    Since \(0\le \ell_u(A_u)\le1\), the contraction inequality used in
    \cref{proof:appendix_f_one_reduction} gives
    \begin{align}
        \phi_\alpha(W)
        =
        \phi_\alpha\left(\ell_u(A_u)(\rho-1)\right)
        \le
        \phi_\alpha(\rho-1)
        =
        f_\alpha(\rho).
        \label{eq:phi_contraction_for_pointwise_martingale}
    \end{align}
    Therefore,
    \begin{align}
        \E_{u-1}\left[
            \exp\left(
            -\eta(\hf_u(z)-\bar{\ell}_u(z))
            \right)
            \right]
         & =
        \E_{u-1}\left[
            \exp\left(
            -\alpha(W-\E_{u-1}[W])
            \right)
            \right]
        \\
         & =
        e^{\alpha \E_{u-1}[W]}
        \left(
        1-\alpha\E_{u-1}[W]
        +
        \E_{u-1}[\phi_\alpha(W)]
        \right)
        \\
         & \le
        e^{\alpha \E_{u-1}[W]}
        \left(
        1-\alpha\E_{u-1}[W]
        +
        \E_{u-1}[f_\alpha(\rho)]
        \right)
        \\
         & \le
        \exp\left(
        \E_{u-1}[f_\alpha(\rho)]
        \right),
    \end{align}
    where the last step uses \(1+x\le e^x\). This proves the supermartingale property by iteration.
    The high-probability statement follows from Ville's inequality.
\end{proof}

Combining \eqref{eq:martingale_weight_decomposition_bar_ell} with
\eqref{eq:pointwise_exponent_bound_no_alpha_loss}, for a fixed \(z\), with probability at least
\(1-\delta\), simultaneously over \(s\),
\begin{align}
    \exp\left(
    -\eta\sum_{u=1}^s \hf_u(z)
    \right)
    \le
    \exp\left(
    -\eta\sum_{u=1}^s \bar{\ell}_u(z)
    \right)
    \exp\left(
    \Gamma_s(z)+\log\frac1\delta
    \right).
    \label{eq:actual_weight_controlled_by_predictable_weight}
\end{align}
This is the useful form of the statement. It should be read as a bound on the random
\emph{multiplicative change of the density}, not primarily as a bound on the cumulative estimator
itself. If one divides \eqref{eq:pointwise_exponent_bound_no_alpha_loss} by \(\eta\), or equivalently
works with the unscaled variable
\[
    \ell_u(A_u)(r_u(A_u,z)-1),
\]
then one obtains a term of order \(\eta^{-1}\log(1/\delta)\), or \(\alpha^{-1}\log(1/\delta)\).
That form is usually too crude to be the final high-probability regret argument.

The point of the proposition is instead conceptual and structural: the same quantity that controls
the variance term also controls the upward spikes of the exponential-weights density. Indeed, if
\(T_0^{(u)}\sim q_u\) is independent of \(A_u\sim p_u\), then averaging the one-step budget over
\(z\sim q_u\) gives
\begin{align}
    \E_{T_0^{(u)}\sim q_u}
    \E_{A_u\sim p_u}
    \left[
    f_\alpha(r_u(A_u,T_0^{(u)}))
    \right]
    =
    D_{f_\alpha}
    \left(
    \PP_{A_u,T_0^{(u)}}
    \middle\|
    \PP_{A_u}\otimes \PP_{T_0^{(u)}}
    \right)
    =
    I_{f_\alpha}(A_u;T_0^{(u)}).
    \label{eq:pointwise_budget_averages_to_variance}
\end{align}
Thus the \(f_\alpha\)-divergence has two interpretations. Averaged over \(z\sim q_u\), it is exactly
the variance term in the expected regret bound. Pointwise in \(z\), it is the exponential compensator
for negative estimator fluctuations, which are precisely the fluctuations that make the density at
\(z\) grow.